\chapter{Bedwetting (Nocturnal Enuresis)}

\section*{Definition}
Bedwetting (nocturnal enuresis) is involuntary urination during sleep in a child aged five years or older. It affects approximately 15-20 percent of five-year-olds and resolves spontaneously in most cases, with an annual remission rate of about 15 percent.

\section*{When to See a Doctor}
\begin{itemize}
\item Bedwetting in a child aged five years or older that concerns the child or family
\item Daytime wetting (suggests more complex voiding dysfunction)
\item Painful urination, excessive thirst, or weight loss
\item Constipation or soiling (encopresis)
\item Snoring or breathing difficulties during sleep
\item New-onset bedwetting after six months of dryness (secondary enuresis)
\item History of urinary tract infections
\end{itemize}

\section*{How It Develops}
Bedwetting involves a mismatch between your bladder's capacity to store urine overnight and your body's ability to either produce less urine or awaken in response to a full bladder. Three main mechanisms contribute: (1) nocturnal polyuria from insufficient antidiuretic hormone (ADH) secretion at night; (2) reduced functional bladder capacity; (3) impaired arousal from deep sleep in response to bladder distention.

\section*{Causes}
\begin{itemize}
\item Delayed maturation of the neural control of urination
\item Insufficient nocturnal ADH secretion (nocturnal polyuria)
\item Small functional bladder capacity
\item Deep sleep pattern with impaired arousal
\item Genetic predisposition (family history)
\item Constipation pressing on the bladder
\item Burning urination, fever, or other symptoms suggesting a urinary tract infection
\item Type 1 diabetes (polydipsia and polyuria)
\item Psychological stress
\item Obstructive sleep apnea
\end{itemize}

\section*{Related Symptoms}
\begin{itemize}
\item Frequent urination (urinary frequency)
\item Urinary urgency
\item Urinary tract infection
\item Constipation
\item Daytime wetting (diurnal enuresis)
\item Sleep apnea or snoring
\item Nocturnal polyuria
\item Excessive thirst, large amounts of urine, or weight loss
\end{itemize}

\section*{Medical Evaluation}
\begin{itemize}
\item Complete history and voiding diary
\item Urinalysis is not routinely needed for typical primary bedwetting, but is considered for recent onset, daytime symptoms, illness, or symptoms suggesting infection or diabetes
\item Assessment of bowel function (constipation evaluation)
\item Bladder ultrasound or flow testing only if daytime voiding dysfunction, recurrent infections, abnormal examination, or poor response suggests a specific urinary problem
\item Sleep study if obstructive sleep apnea is suspected
\end{itemize}

\section*{Treatment Options}
\begin{itemize}
\item Address fluid intake, regular daytime toileting, constipation, and the child’s concerns first; then offer a moisture alarm when appropriate. Alarms generally provide better long-term benefit than medication.
\item Desmopressin may be prescribed, usually for children older than 7 years, when rapid or short-term dryness is important or an alarm is unsuitable. Strict fluid restriction around each dose is essential.
\item Oxybutynin or another bladder medicine is not routine treatment for isolated bedwetting and should be used only under specialist guidance, usually when daytime bladder symptoms are present.
\item Treat constipation with an age-appropriate bowel plan advised by a clinician
\item Address any underlying psychological or sleep disorder
\item Combination therapy or specialist referral may be considered after recurrence or failure of appropriate alarm and/or desmopressin treatment
\item Reassurance and patience, as most cases resolve spontaneously
\end{itemize}

\section*{Self-Care and Home Management}
\begin{itemize}
\item Offer adequate fluids throughout the day and avoid excessive drinks close to bedtime; do not make the child dehydrated.
\item Encourage regular toileting before sleep (double voiding)
\item Use a moisture alarm to condition arousal to bladder fullness
\item Protect the mattress with a waterproof cover
\item Praise helpful behaviors, such as using the toilet before bed and participating in the plan; never punish or shame wet nights.
\item Ensure regular bowel movements to prevent constipation
\item Avoid caffeine and consider limiting fizzy drinks if they worsen symptoms.
\end{itemize}

\section*{References}
\begin{thebibliography}{99}
\bibitem{bedwetting:ref1} Nevéus T, Fonseca E, et al. ``Management of nocturnal enuresis in children.'' \textit{Lancet}. 2020;396(10250):643-654.
\bibitem{bedwetting:ref2} Kiddoo DA. ``Nocturnal enuresis.'' \textit{BMJ Clin Evid}. 2015;2015:0305.
\bibitem{bedwetting:ref3} Vande Walle J, Rittig S, et al. ``Practical consensus guidelines for the management of enuresis.'' \textit{Eur J Pediatr}. 2012;171(6):971-983.
\bibitem{bedwetting:ref4} Kliegman RM, St Geme JW, et al. \textit{Nelson Textbook of Pediatrics}, 21st ed. Elsevier, 2020.
\bibitem{bedwetting:ref5} National Institute for Health and Care Excellence. ``Bedwetting in under 19s (CG111).'' Updated 2023. https://www.nice.org.uk/guidance/cg111.
\end{thebibliography}
